\chapter{Trading de l'énergie et du gaz}\label{ch:s2:power-and-gas-trading}

L'électricité achetée la veille pour livraison le lendemain et l'électricité achetée une heure avant la livraison sont deux produits différents. L'enchère day-ahead se fixe
sur les prévisions éoliennes et solaires de la veille au matin~; le marché infrajournalier valorise ce qui souffle et brille réellement. Kiesel et Paraschiv ont montré que les prix infrajournaliers allemands
s'ajustent aux erreurs de prévision des énergies renouvelables. Sur les prix day-ahead allemands de 2024--2025, avec des prix infrajournaliers simulés à partir d'erreurs de prévision
éoliennes introduites volontairement, une batterie de 1~MW, 2~MWh programmée sur l'enchère day-ahead rapporte 69\,900~\euro par an. Reprogrammer la batterie en fonction des prix infrajournaliers ajoute
2\,900~\euro. Un trader qui prévoit la moitié de l'erreur éolienne gagne 2.54~\euro par MWh échangé entre les deux marchés. La construction est \texttt{firm.powergas}.

\section{Day-ahead contre infrajournalier}

\begin{definition}[Spread day-ahead--infrajournalier]\label{def:s2:power-and-gas-trading:daid}
Le \emph{spread day-ahead--infrajournalier}\index{spread day-ahead--infrajournalier} pour une heure de livraison est le prix de l'enchère day-ahead diminué du prix infrajournalier pour
la même heure~; il est piloté par ce qui a changé entre l'enchère et la livraison, surtout les prévisions d'éolien, de solaire et de demande.
\end{definition}

Les prix day-ahead allemands sont étroitement liés à la charge résiduelle, c'est-à-dire la charge qui reste après l'éolien et le solaire. Sur 2024 et 2025, une régression du prix horaire sur
la charge résiduelle donne une pente de 3.05~\euro par MWh et par GW et une corrélation de 0.83. Si l'enchère s'est fixée sur une prévision éolienne inférieure de 1~GW à la réalité, la
tranche horaire a été valorisée pour 1~GW de charge résiduelle de plus qu'il ne s'en est présenté, et le prix infrajournalier devrait être environ 3~\euro plus bas. \texttt{firm.powergas} construit les prix infrajournaliers
de cette manière~: le prix day-ahead diminué de 3.05~\euro par GW d'éolien imprévu, avec des erreurs de prévision d'écart-type 2~GW qui persistent d'une heure à
l'autre (autocorrélation 0.9), plus 6~\euro d'autre bruit (\cref{lst:s2:power-and-gas-trading:trade}). Ces erreurs sont introduites volontairement~: les prix infrajournaliers réels et les prévisions
ne sont pas libres d'accès.

Les résultats empiriques sur les prix infrajournaliers réels sont contrastés dans le détail. Kiesel et Paraschiv ont montré que les prix infrajournaliers s'ajustent de façon asymétrique aux erreurs de prévision renouvelables
et au volume échangé, selon un seuil dans la demande censée être couverte par la capacité conventionnelle programmée. Ziel, avec un
modèle de régression sur les prix infrajournaliers allemands, n'a trouvé aucune preuve statistiquement significative que les erreurs éoliennes ou solaires positives et négatives aient
des effets différents.

Un trader dont la prévision de l'erreur a une corrélation de 0.5 avec elle vend sur le day-ahead et rachète en infrajournalier lorsqu'il anticipe davantage de vent que
l'enchère n'en supposait, et fait l'inverse lorsqu'il en anticipe moins~; avec une demi-fourchette de 1~\euro du côté infrajournalier~:

\begin{center}
\small
\begin{tabular}{@{}lccc@{}}
\toprule
habileté de prévision (corrélation), demi-fourchette & \euro{} par MWh échangé & \euro{} par an et par MW & IC\\
\midrule
0.5, 1~\euro & 2.54 & 13\,600 & 0.36\\
0.3, 1~\euro & 1.08 & 5\,800 & 0.21\\
0.5, 3~\euro & 0.54 & 2\,900 & 0.36\\
\bottomrule
\end{tabular}
\end{center}

La stratégie est présente sur le marché dans 61~\% des heures. L'edge est une prévision de l'erreur, et il diminue vite avec l'habileté et avec les coûts.

\section{Positions météo}

\begin{definition}[Position météo]\label{def:s2:power-and-gas-trading:weather}
Une \emph{position météo}\index{position météo} est une position sur l'électricité ou le gaz prise sur une prévision météorologique (température, vent,
nébulosité, précipitations) qui diffère de celle que reflètent les prix de marché, et clôturée lorsque les prévisions convergent ou que le temps se présente.
\end{definition}

La stratégie day-ahead est la \omterm{def:s2:power-and-gas-trading:weather}{position météo} la plus courte. Sur plusieurs jours et plusieurs semaines, les prévisions de température gouvernent la demande de chauffage et de climatisation, et avec
elle les prix du gaz et de l'électricité. Un trader qui exploite une révision de prévision avant que le marché ne l'intègre dans ses prix dispose du même type d'avantage, sur un horizon plus long et à
une plus grande échelle, avec les mêmes ennemis~: une prévision que le marché possède déjà, et les coûts. Le livre~3 (chapitre~11) traite des dérivés climatiques, qui
couvrent la météo elle-même.

\section{Batteries et actifs flexibles}

\begin{definition}[Arbitrage par batteries]\label{def:s2:power-and-gas-trading:battery}
L'\emph{arbitrage par batteries}\index{arbitrage par batteries} consiste à programmer une batterie pour charger aux heures bon marché et décharger aux heures chères, d'abord
contre l'enchère day-ahead, puis en reprogrammant en fonction des prix infrajournaliers, dans les limites de sa puissance, de son énergie, de son rendement et de son nombre de cycles.
\end{definition}

Les prix allemands ont la forme dont une batterie a besoin (\cref{fig:s2:power-and-gas-trading:profile})~: bon marché à midi, quand le solaire culmine, chers le soir.
Les prix ont été négatifs pendant 457 heures en 2024 et 573 en 2025, jusqu'à $-$250.32~\euro~; à 13~h, heure locale, 23.4~\% des heures étaient négatives, et à 19~h
aucune. La batterie est programmée chaque jour par programmation dynamique sur son état de charge et sur l'énergie chargée jusque-là
(\cref{lst:s2:power-and-gas-trading:battery}), avec un rendement aller-retour de 88~\% et un cycle complet par jour. Sur les 727 jours de 24 heures, en connaissant
à l'avance les prix de l'enchère de chaque jour, comme dans le livre~3, elle rapporte 69\,900~\euro par MW et par an. Un quart de ses heures de charge (24.8~\%) ont des prix négatifs.

\begin{omfigure}
\begin{tikzpicture}
\begin{axis}[width=11cm,height=5.2cm,axis y line*=right,axis x line=none,ylabel={\small heures négatives (\%)},
  tick label style={font=\footnotesize},xmin=-0.5,xmax=23.5,ymin=0,ymax=60,ybar,bar width=5pt,table/col sep=comma]
\addplot[fill=gray!30,draw=gray] table[x=hour,y=negative]{figdata/strategies-2/23-power-and-gas-trading/profile.csv};\label{plot:s2:powneg}
\end{axis}
\begin{axis}[width=11cm,height=5.2cm,axis y line*=left,xlabel={\small heure, heure locale allemande},ylabel={\small \euro/MWh, moyenne},
  tick label style={font=\footnotesize},xmin=-0.5,xmax=23.5,ymin=0,ymax=150,xtick={0,3,6,9,12,15,18,21},grid=major,grid style={gray!25},
  table/col sep=comma,legend style={font=\scriptsize,at={(0.5,-0.3)},anchor=north},legend columns=3]
\addplot[omThm,thick,mark=*,mark size=1pt] table[x=hour,y=p2024]{figdata/strategies-2/23-power-and-gas-trading/profile.csv};
\addlegendentry{2024}
\addplot[omDef,thick,mark=square*,mark size=1pt] table[x=hour,y=p2025]{figdata/strategies-2/23-power-and-gas-trading/profile.csv};
\addlegendentry{2025}
\addlegendimage{/pgfplots/refstyle=plot:s2:powneg}\addlegendentry{heures négatives (droite)}
\end{axis}
\end{tikzpicture}

\omcaption{Prix day-ahead allemands (DE-LU) par heure de la journée, moyennés sur 2024 et sur 2025, et part des heures à prix négatifs sur les deux
années. Données~: Bundesnetzagentur / SMARD.de (CC BY 4.0).}\label{fig:s2:power-and-gas-trading:profile}
\end{omfigure}

La batterie qui a vendu son programme day-ahead peut le modifier en infrajournalier~: racheter l'électricité qu'elle ne veut plus livrer, ou en vendre davantage, en payant la
demi-fourchette sur chaque MWh modifié. En reprogrammant en fonction des prix infrajournaliers de la journée, le programme change dans 71~\% des jours~:

\begin{center}
\small
\begin{tabular}{@{}p{5.0cm}>{\centering\arraybackslash}p{4.2cm}>{\centering\arraybackslash}p{4.2cm}@{}}
\toprule
1~MW, 2~MWh, rendement aller-retour de 88~\% & day-ahead (\euro{} par an) & supplément infrajournalier (\euro{} par an)\\
\midrule
un cycle par jour, demi-fourchette 1~\euro & 69\,900 & 2\,900\\
un cycle, sans spread & 69\,900 & 3\,800\\
un cycle, demi-fourchette 3~\euro & 69\,900 & 1\,600\\
deux cycles par jour, demi-fourchette 1~\euro & 83\,800 & 4\,900\\
\bottomrule
\end{tabular}
\end{center}

Les deux colonnes sont généreuses. La colonne day-ahead connaît les prix de l'enchère avant d'y soumettre ses offres~; la colonne infrajournalière reprogramme une seule fois en connaissant
l'ensemble des prix infrajournaliers de la journée. Le supplément infrajournalier est faible parce que les erreurs introduites volontairement ne déplacent les prix que de quelques euros, alors que l'écart de prix de la journée
moyenné sur une centaine (\cref{fig:s2:power-and-gas-trading:battery}).

\begin{omfigure}
\begin{tikzpicture}
\begin{axis}[width=11cm,height=5cm,xlabel={\small jour (1 janvier 2024 au 31 décembre 2025, jours de changement d'heure omis)},
  ylabel={\small cumul en milliers d'\euro{} par MW},tick label style={font=\footnotesize},xmin=0,xmax=730,ymin=0,ymax=155,
  grid=major,grid style={gray!25},table/col sep=comma,legend style={font=\scriptsize,at={(0.03,0.97)},anchor=north west}]
\addplot[omThm,thick] table[x=day,y=da]{figdata/strategies-2/23-power-and-gas-trading/battery.csv};
\addplot[omDef,thick,dashed] table[x=day,y=total]{figdata/strategies-2/23-power-and-gas-trading/battery.csv};
\legend{programme day-ahead,avec réoptimisation infrajournalière}
\end{axis}
\end{tikzpicture}

\omcaption{Revenu cumulé d'une batterie de 1~MW, 2~MWh sur les prix day-ahead allemands, 2024--2025, et avec son programme réoptimisé face aux
prix infrajournaliers synthétiques. Données~: Bundesnetzagentur / SMARD.de (CC BY 4.0)~; \texttt{s2\_powergas.battery\_table}.}\label{fig:s2:power-and-gas-trading:battery}
\end{omfigure}

\section{Risque dans les portefeuilles d'électricité}

L'électricité ne se stocke pas en masse, si bien que chaque heure constitue un produit à part entière, et un portefeuille d'électricité est un portefeuille de 8\,760 contrats par an, chacun soumis à sa propre météo.
Ses risques sont les erreurs de prévision elles-mêmes, qui sont corrélées d'une heure à l'autre (une mauvaise prévision éolienne se trompe toute l'après-midi)~; des prix qui bondissent
de plusieurs centaines d'euros lorsque quelques gigawatts font défaut (le maximum de 2024 a atteint 936.28~\euro)~; les pénalités de déséquilibre pour une livraison inférieure à ce qui a été vendu (livre~3,
chapitre~6)~; et les limites de l'actif physique, à savoir une batterie pleine au moment où les prix deviennent négatifs. Une batterie se dégrade aussi à chaque cycle, ce que le
tableau ci-dessus ne facture pas. Un second cycle par jour ajoute 13\,800~\euro de revenus day-ahead, au prix d'une durée de vie de la batterie que son propriétaire doit valoriser.

\section{Fiches de stratégie}

\begin{strategyfile}{Transaction infrajournalière sur erreur de prévision}\label{strat:s2:power-and-gas-trading:forecast}
\sfield{Qui vous paie, et pourquoi} Les producteurs et fournisseurs qui doivent équilibrer leurs erreurs de prévision en infrajournalier.
\sfield{Instruments et plateformes} Enchère day-ahead~; marché infrajournalier continu.
\sfield{Signal} La prévision du trader pour l'éolien, le solaire et la charge, comparée à celle de l'enchère.
\sfield{Dimensionnement et exécution} Positions proportionnées à la confiance dans la prévision~; clôturées en infrajournalier.
\sfield{Coûts} Fourchette acheteur-vendeur infrajournalière~; pénalités de déséquilibre sur ce qui n'est pas clôturé.
\sfield{Comment elle meurt} Prévisions publiques meilleures~; coûts.
\sfield{Horizon, capacité, infrastructure} Heures~; modèles météorologiques et accès infrajournalier rapide.
\sfield{Backtest honnête} Prévisions à la date de l'enchère, et non postérieures.
\sfield{Sources} Kiesel et Paraschiv (2017)~; Ziel (2017)~; ce chapitre~: 2.54~\euro par MWh pour une habileté de 0.5.
\end{strategyfile}

\begin{strategyfile}{Position météo sur le gaz}\label{strat:s2:power-and-gas-trading:weather}
\sfield{Qui vous paie, et pourquoi} Les acteurs de marché qui valorisent la prévision de la veille.
\sfield{Instruments et plateformes} Contrats à terme gaz pour les semaines à venir~; base électrique.
\sfield{Signal} Révisions des prévisions de degrés-jours de chauffage et de climatisation, comparées à la prévision déjà cotée.
\sfield{Dimensionnement et exécution} Proportionnellement à l'effet attendu de la révision sur la demande.
\sfield{Coûts} Fourchette acheteur-vendeur des contrats à terme.
\sfield{Comment elle meurt} Tout le monde dispose des mêmes modèles~; les révisions sont cotées en quelques minutes.
\sfield{Horizon, capacité, infrastructure} De quelques jours à quelques semaines.
\sfield{Backtest honnête} Versions successives des prévisions telles qu'elles ont été publiées.
\sfield{Sources} Aucun chiffre de performance vérifié.
\end{strategyfile}

\begin{strategyfile}{Arbitrage par batteries en day-ahead}\label{strat:s2:power-and-gas-trading:battery}
\sfield{Qui vous paie, et pourquoi} La forme du prix journalier~: le solaire à midi, la demande le soir.
\sfield{Instruments et plateformes} Une batterie~; enchère day-ahead~; marché infrajournalier~; réserves.
\sfield{Signal} Prévoir les prix day-ahead, puis les prix infrajournaliers.
\sfield{Dimensionnement et exécution} Un programme journalier par programmation dynamique~; réoptimisé en infrajournalier.
\sfield{Coûts} Pertes aller-retour~; dégradation~; spreads.
\sfield{Comment elle meurt} D'autres batteries aplatissant la forme.
\sfield{Horizon, capacité, infrastructure} Heures~; l'actif.
\sfield{Backtest honnête} Prévisions de prix, et non prix réalisés aux enchères~; dégradation par cycle.
\sfield{Sources} Ce chapitre~: 69\,900~\euro par MW et par an en day-ahead avec connaissance du futur, 2\,900~\euro de plus en infrajournalier.
\end{strategyfile}

\begin{strategyfile}{Optimisation d'une centrale à spark spread}\label{strat:s2:power-and-gas-trading:plant}
\sfield{Qui vous paie, et pourquoi} Le spread entre l'électricité et le combustible durant les heures où une centrale fonctionne.
\sfield{Instruments et plateformes} Une centrale à gaz ou un contrat de tolling~; contrats à terme sur l'électricité, le gaz et le carbone.
\sfield{Signal} Le clean spark spread heure par heure, comparé aux coûts de la centrale.
\sfield{Dimensionnement et exécution} Mise en service lorsque le spread dépasse les coûts~; couverture à terme par blocs.
\sfield{Coûts} Coûts de démarrage~; durées minimales de fonctionnement.
\sfield{Comment elle meurt} Les énergies renouvelables qui prennent les heures de la centrale.
\sfield{Horizon, capacité, infrastructure} De quelques heures à quelques années.
\sfield{Backtest honnête} Contraintes de démarrage et de montée en charge.
\sfield{Sources} Aucun chiffre de performance vérifié.
\end{strategyfile}

\begin{strategyfile}{Heures à prix négatifs}\label{strat:s2:power-and-gas-trading:negative}
\sfield{Qui vous paie, et pourquoi} Les producteurs qui doivent ou veulent produire à n'importe quel prix (mécanismes de subvention, centrales peu flexibles).
\sfield{Instruments et plateformes} Charge flexible ou stockage~; marchés day-ahead et infrajournalier.
\sfield{Signal} Prévision du solaire et de l'éolien comparée à la charge à midi et les jours fériés.
\sfield{Dimensionnement et exécution} Consommer ou charger pendant les heures négatives.
\sfield{Coûts} Frais de réseau sur la consommation.
\sfield{Comment elle meurt} Des règles de subvention qui cessent de payer pendant les heures négatives~; une demande plus flexible.
\sfield{Horizon, capacité, infrastructure} Heures.
\sfield{Backtest honnête} Frais de réseau~; la prévision, et non le prix réalisé.
\sfield{Sources} Données SMARD~: 457 heures négatives en 2024 et 573 en 2025.
\end{strategyfile}

\section{Tutoriel~: la prévision éolienne}\label{tut:s2:power-and-gas-trading:tut}

\begin{tutorial}
\textbf{Objectif.} Mesurer le profil des prix allemands et les heures négatives, simuler des prix infrajournaliers à partir des erreurs de prévision éolienne, négocier les erreurs, et programmer une batterie
day-ahead puis en infrajournalier. \textbf{État final~:} les trois tableaux et les deux figures.
\begin{tutsteps}
\item \textbf{Prix infrajournaliers et transaction sur prévision}.
\omcode{code/firm/powergas/firm_powergas.py}{42}{63}{Les erreurs éoliennes déplacent les prix infrajournaliers~; une prévision partielle permet de les négocier.}\label{lst:s2:power-and-gas-trading:trade}
\item \textbf{La batterie}.
\omcode{code/firm/powergas/firm_powergas.py}{66}{101}{Programmation dynamique journalière, en day-ahead ou face à un programme engagé.}\label{lst:s2:power-and-gas-trading:battery}
\item \textbf{Exécuter} \texttt{real()}, \texttt{trade()}, \texttt{battery\_table()} et \texttt{fig\_powergas.py}.
\end{tutsteps}
\textbf{Pour aller plus loin.} Engager la batterie sur une prévision de prix plutôt que sur les prix de l'enchère~; réoptimiser en infrajournalier heure par heure~; ajouter
un coût de dégradation par cycle.
\end{tutorial}

\section{Construction~: électricité et gaz}\label{bld:s2:power-and-gas-trading:powergas}

\begin{build}
\bfield{Rôle} Prix infrajournaliers à partir des erreurs de prévision, transaction sur erreur de prévision, et programmation de la batterie en day-ahead et en infrajournalier.
\bfield{Interface} \texttt{PowerConfig(\dots)}, \texttt{intraday(da, cfg)}, \texttt{forecast\_trade(da, idp, err, cfg)},
\texttt{battery(prices, cfg, committed)}.
\bfield{Règles} Charger coûte $1/\sqrt{\eta}$ de l'énergie du réseau, décharger restitue $\sqrt{\eta}$~; la batterie se vide chaque jour~; les écarts
par rapport à un programme engagé coûtent la demi-fourchette.
\bfield{Tests d'acceptation} \texttt{code/firm/powergas/tests/}~: pas d'erreur, pas de spread~; une prévision parfaite ne perd jamais sans coûts~; la batterie calculée à la main~;
réoptimiser sur les prix engagés ne change rien.
\bfield{Pour aller plus loin} Quarts d'heure comme produits~; réserves~; dégradation.
\end{build}

\begin{omsources}
\item R.~Kiesel and F.~Paraschiv, «~Econometric analysis of 15-minute intraday electricity prices~», \emph{Energy Economics} 64, 2017.
\item F.~Ziel, «~Modeling the impact of wind and solar power forecasting errors on intraday electricity prices~», \emph{14th International Conference
on the European Energy Market}, 2017.
\item Bundesnetzagentur / SMARD.de, prix J-1 horaires, production et charge pour l'Allemagne, 2024--2025, CC BY 4.0 (fichier du livre~3).
\end{omsources}

\section{Exercices}

\begin{exercise}[$\star$]\label{exo:s2:power-and-gas-trading:1}
Le vent dépasse de 1.5~GW la prévision de l'enchère. À 3.05~\euro par MWh et par GW, de combien le prix infrajournalier doit-il se situer sous le prix day-ahead~?
\end{exercise}

\begin{exercise}[$\star$]\label{exo:s2:power-and-gas-trading:2}
Une batterie achète 2~MWh de stockage à 0~\euro et le revend à 100~\euro avec un rendement aller-retour de 88~\%. Que gagne-t-elle~?
\end{exercise}

\begin{exercise}[$\star$]\label{exo:s2:power-and-gas-trading:3}
La transaction sur erreur de prévision rapporte 2.54~\euro par MWh sur 61~\% des heures, avec 1~MW. Vérifier le chiffre annuel sur deux années de 727 jours.
\end{exercise}

\begin{exercise}[$\star\star$]\label{exo:s2:power-and-gas-trading:4}
Pourquoi l'edge de la transaction sur erreur de prévision chute-t-il si vite avec une demi-fourchette de 3~\euro~?
\end{exercise}

\begin{exercise}[$\star\star$]\label{exo:s2:power-and-gas-trading:5}
Pourquoi les prix allemands sont-ils le plus souvent négatifs à midi~?
\end{exercise}

\begin{exercise}[$\star\star$]\label{exo:s2:power-and-gas-trading:6}
Kiesel et Paraschiv ont trouvé une réponse asymétrique aux erreurs de prévision~; Ziel n'en a trouvé aucune significative. Comment les deux peuvent-ils avoir raison~?
\end{exercise}

\begin{exercise}[$\star\star\star$]\label{exo:s2:power-and-gas-trading:7}
\emph{Programmation.} Exécuter \texttt{battery\_table(PowerConfig(cycles=2))}. Qu'apporte le deuxième cycle, et que laisse de côté le tableau~?
\end{exercise}

\begin{exercise}[$\star\star\star$]\label{exo:s2:power-and-gas-trading:8}
\emph{Trouver la faille.} «~Notre backtest de batterie a rapporté 69\,900~\euro par MW et par an sur les prix day-ahead allemands~; c'est ce que nous gagnerons.~»
\end{exercise}

\section{Problème~: La prévision éolienne}

\begin{problem}[{Problème du week-end --- trading d'électricité et de gaz}]\label{pb:s2:power-and-gas-trading:1}
Les prix day-ahead de l'Allemagne, le marché infrajournalier synthétique et les données publiques.

\medskip\noindent\textbf{Partie I --- Les marchés.}
\begin{enumerate}
\item Définir le \omterm{def:s2:power-and-gas-trading:daid}{spread day-ahead--infrajournalier}.
\item Donner la pente et la corrélation du prix en fonction de la charge résiduelle.
\item Comment les prix infrajournaliers synthétiques sont-ils construits, et quel élément y est implanté~?
\item Qu'ont trouvé Kiesel et Paraschiv et Ziel~?
\end{enumerate}

\medskip\noindent\textbf{Partie II --- La transaction sur erreur de prévision.}
\begin{enumerate}[resume]
\item Décrire la transaction.
\item Donner ses résultats pour chaque niveau d'habileté et chaque coût.
\item Définir une \omterm{def:s2:power-and-gas-trading:weather}{position météo}.
\item Comment la transaction s'éteint-elle~?
\end{enumerate}

\medskip\noindent\textbf{Partie III --- La batterie.}
\begin{enumerate}[resume]
\item Définir l'\omterm{def:s2:power-and-gas-trading:battery}{arbitrage par batteries}.
\item Décrire le profil de prix et les heures négatives.
\item Comment la batterie est-elle programmée~?
\item Donner le tableau de la batterie.
\end{enumerate}

\medskip\noindent\textbf{Partie IV --- Le verdict.}
\begin{enumerate}[resume]
\item Énoncer le \emph{résultat nommé}~: le revenu de la batterie en day-ahead seul et avec réoptimisation infrajournalière.
\item Pourquoi les deux colonnes sont-elles généreuses~?
\item Pourquoi le supplément infrajournalier est-il faible ici~?
\item Quels sont les risques d'un book d'électricité~?
\item Comment backtester honnêtement la batterie~?
\item Quelle fiche de stratégie meurt en premier à mesure que les batteries se multiplient~?
\item Quel est le lien entre la transaction sur erreur de prévision et les transactions de flux du chapitre~15~?
\item En une phrase~: que vend un trader d'électricité~?
\end{enumerate}
\end{problem}

\section{Questions d'entretien}

\begin{interviewq}[$\star$ \iqroles{trader}]\label{iq:s2:power-and-gas-trading:1}
Pourquoi les prix de l'électricité peuvent-ils être négatifs~?
\end{interviewq}

\begin{interviewq}[$\star\star$ \iqroles{chercheur}]\label{iq:s2:power-and-gas-trading:2}
Comment prévoiriez-vous le \omterm{def:s2:power-and-gas-trading:daid}{spread day-ahead--infrajournalier} pour les heures de demain~?
\end{interviewq}

\begin{interviewq}[$\star\star$ \iqroles{trader}]\label{iq:s2:power-and-gas-trading:3}
Votre prévision éolienne dépasse de 3~GW celle du marché pour l'après-midi de demain. Que faites-vous, et pour quel volume~?
\end{interviewq}

\begin{interviewq}[$\star\star$ \iqroles{risque}]\label{iq:s2:power-and-gas-trading:4}
Comment fixeriez-vous des limites pour un desk d'électricité infrajournalier~?
\end{interviewq}

\begin{interviewq}[$\star\star$ \iqroles{développeur}]\label{iq:s2:power-and-gas-trading:5}
Concevez le système qui programme une flotte de batteries à chaque quart d'heure.
\end{interviewq}

\begin{interviewq}[$\star\star\star$ \iqroles{chercheur}]\label{iq:s2:power-and-gas-trading:6}
Une batterie de rendement aller-retour $\eta$ achète à $p_1$ et vend à $p_2$. Montrez que la transaction ne rapporte que si $p_2 > p_1/\eta$ (pour des
prix positifs), et dites ce qui change lorsque $p_1 < 0$.
\end{interviewq}
